% !TEX root = ../main.tex

\section{Conclusiones y trabajo futuro}
\label{sec:conclusion}

Este art\'iculo present\'o \sculpt, un lenguaje de programaci\'on
tangible cuyas piezas se imprimen en 3D y se ensamblan f\'isicamente
para construir esculturas Turing-completas. Sobre un conjunto m\'inimo
de 13 instrucciones que operan sobre pilas FILO, se demostr\'o
formalmente la completitud del lenguaje mediante una traducci\'on
$\Phi$ que mapea m\'aquinas de Turing arbitrarias a programas \sculpt.
La libertad est\'etica de las piezas distingue al lenguaje frente a las tres
familias de trabajo relacionado y, al mismo tiempo, ampl\'ia las
posibilidades de accesibilidad y de uso expresivo.

La validaci\'on con participantes mostr\'o un patr\'on consistente.
La fisicalidad funciona como puerta de entrada y los componentes son
reconocibles para audiencias diversas, pero programar sigue requiriendo
razonar algor\'itmicamente y eso no se elimina por cambiar el soporte.
La percepci\'on docente, m\'as positiva que la de los aprendices,
sugiere que el potencial pedag\'ogico de \sculpt es mayor que el que
sus resultados de corto plazo permiten medir. Los artistas, por su
parte, encontraron en el lenguaje un medio expresivo en s\'i mismo,
m\'as all\'a de su capacidad computacional sin dejarla de lado.

\paragraph{Trabajo futuro.}
Quedan abiertas varias l\'ineas. Primero, estudios m\'as extensos en
contextos escolares y talleres, con cohortes estables que permitan
medir el aprendizaje longitudinal y no solo el desempe\~no en sesi\'on
\'unica. Segundo, mejoras a \sculpter como herramienta de desarrollo y como gui\'a de referencia.
Agregar elementos como un sistema de manejo de errores m\'as informativo, \emph{language server} con
\emph{syntax highlighting} y completaci\'on, y un visualizador que
permita ver el ensamble final antes de imprimir las piezas. Tercero,
herramientas dirigidas espec\'ificamente a artistas. Herramientas como exportaci\'on de
los bloques de una escultura como modelos 3D listos para impresi\'on,
y un \emph{esc\'aner de esculturas} capaz de leer una construcci\'on
f\'isica y reconstruir el programa correspondiente por visi\'on por
computador. Cuarto, exploraci\'on de \sculpt en contextos art\'isticos
ampliados como instalaciones interactivas, performances, talleres
multidisciplinarios.

Esta \'ultima l\'inea el lenguaje como objeto de \emph{programming
experience}, donde lo que importa no es solo qu\'e se computa sino
qu\'e se siente al hacerlo es la que m\'as promete. 
La fisicalidad y abstracci\'on geometrica de las sentencias de \sculpt
es m\'as que una particularidad es el llamado
a repensar qu\'e es crear codigo, qui\'en puede hacerlo y por qu\'e lo har\'ia.
El formalismo y descripci\'on l\'ogica de sus piezas es la invitaci\'on a replantear
que es crear con c\'odigo, qu\'ien puede hacerlo y por qu\'e lo har\'ia.
El c\'odigo como arte, el arte como c\'odigo.

\endinput
